\documentclass[12pt]{article}
\usepackage[utf8]{inputenc}
\usepackage[T1]{fontenc}
\usepackage{hyperref}
\title{Citak Language Resources for K-12 Curriculum Development}
\author{Compiled by Claude Sonnet 5 (claude-sonnet-5) for Andr\'as Kornai}
\date{}
\begin{document}
\maketitle
\section{Introduction}
Citak (ISO 639-3 txt) is a Trans-New Guinea language of the Citak-Mitak district, Mappi Regency, Papua, Indonesia. It is also called Kaunak, and speakers call it Kau Adagum. The figure of about 8,000 speakers dates from 1985. Searches found almost no usable materials, so this report is brief.
\section{Primary Educational Resources}
None found in search results.
\section{Linguistic and Reference Materials}
\begin{itemize}
\item Citak language, Wikipedia overview: \url{https://en.wikipedia.org/wiki/Citak_language}
\item Glottolog entry for Citak: \url{https://glottolog.org/resource/languoid/id/cita1245.bigmap.html}
\item Trans New Guinea database, Citak wordlist (a few words from van der Wouw 1974): \url{https://transnewguinea.org/language/citak}
\end{itemize}
No grammar was found. Van der Wouw (1974) and Antoni (2010) were mentioned in search results but not accessed, and no URLs were seen for them.
\section{Religious and Community Resources}
A Joshua Project page on the Citak/Tamnim people exists: \url{https://joshuaproject.net/people_groups/11391}. Its content as a language resource was not assessed.
\section{Digital and Audio Resources}
None found in search results.
\section{Curriculum Development Context}
Almost no published description or teaching material is available. Any curriculum would need new documentation work with the community.
\end{document}